% ===== 原文 pp.29--33 =====
% [原文 p.29]
也就是说，一个 $p$ 维算符 $\cO_p$ 在 $G^{(p)}$ 下的荷为
\begin{equation}
  Q(\cO_p)=[\cO_p]\in \mathsf{D}_p\,, \tag{3.7}
\end{equation}
其中 $[\cO_p]$ 是 $\cO_p$ 所在的等价类。

% [Tutor Note: p.29 式 (3.7) 中的 $\mathsf{D}_p$ 就是 (3.4) 定义的“$p$ 维算符按屏蔽关系取等价类”所得集合；此处默认它已具有由 OPE 诱导的交换群结构。核对方法：由 (3.6) $\widehat{G}^{(p)}=\mathsf{D}_p$ 反推，(3.8) 给 $\mathsf{D}_1=\mathsf{D}_{d-3}=\Z$，取 Pontryagin 对偶即得 (3.9) 的两个 $\U(1)$，与 2.2 节从生成元角度得到的结论一致。]

\begin{examplebox}{例 3.1：重访 Maxwell 理论}
我们可以利用上面讨论的屏蔽论证，推断出纯 $\U(1)$ 规范理论具有电与磁的高阶形式对称性。当既无电的物质内容、又无磁的物质内容时，不存在把两个不同 Wilson 算符、或把两个不同 't Hooft 算符联系起来的屏蔽。于是线算符与余维-3 算符的等价类为
\begin{equation}
  \mathsf{D}_1=\Z\,,\qquad \mathsf{D}_{d-3}=\Z\,. \tag{3.8}
\end{equation}
取 Pontryagin 对偶，我们就恢复出此前从另一角度讨论过的电 1-形式对称性与磁 $(d-3)$-形式对称性群
\begin{equation}
  G^{(1)}=U(1)\,,\qquad G^{(d-3)}=U(1) \tag{3.9}
\end{equation}
\end{examplebox}

\subsection{阿贝尔规范理论}

在本小节中，我们使用基于屏蔽的论证来确定具有 阿贝尔规范群的规范理论中出现的高阶形式对称性群。

\subsubsection{含单个物质场的 Maxwell 理论}

\para{设置。} 考虑一个 $d$ 维规范理论，其规范群为 $\U(1)$，并带有一个在 $\U(1)$ 规范群下电荷为 $q\in\Z$ 的物质场 $\phi$\footnote{$\U(1)$ 规范群下的单位电荷，由用于探测该理论的 Wilson 线算符所携带的最小正电荷确定。}。注意，$\phi$ 在时空 Lorentz 群下可以具有任意变换性质（标量、旋量等），并且它可以是无质量场，也可以是有质量场。

我们此前已经看到，对于不含物质场 $\phi$ 的纯 Maxwell 理论，我们有一个 $\U(1)$ 1-形式对称性。现在我们将看到，引入 $\phi$ 会把这一 $\U(1)$ 1-形式对称性破缺为其 $\Z_q$ 子群。也就是说，该理论具有电 1-形式对称性群
\begin{equation}
  G^{(1)}=\Z_q\,. \tag{3.10}
\end{equation}

\para{利用物质场构造非真局域算符。} 关键在于，物质场给出一个非真局域算符 $\phi(x)$，它屏蔽电荷为 $q$ 的 Wilson 线 $\cW_q$，如图 9 所示。

% [原文 p.30]
\begin{figure}[htbp]
\centering
\includegraphics[width=0.28\textwidth]{fig9.png}
\figtitle{9}{电荷为 $q$ 的物质场 $\phi$ 给出一个位于电荷为 $q$ 的 Wilson 线端点的非真局域算符。}
\end{figure}

\begin{definitionbox}{定义 3.2：非真算符}
一个非真 $q$ 维算符是指附着在一族 $p>q$ 维算符上的算符。
\end{definitionbox}

场 $\phi(x)$ 的插入不能定义一个真局域算符，因为这样的算符将不是规范不变的。在规范变换
\begin{equation}
  A(x)\to A(x)-\frac{\dd\theta(x)}{2\pi}\,, \tag{3.11}
\end{equation}
下，插入 $\phi(x)$ 按
\begin{equation}
  \phi(x)\to \ee^{iq\theta(x)}\phi(x) \tag{3.12}
\end{equation}
变换。为使这一插入成为规范不变的，我们可以把它附着到一个电荷为 $q$ 的 Wilson 线算符 $\cW_q(L)$ 上，后者沿一条终止于点 $x$ 的线 $L$ 插入：
\begin{equation}
  \partial L=x\,. \tag{3.13}
\end{equation}
这是因为 $\cW_q(L)$ 的规范变分为
\begin{equation}
  \cW_q(L)=\exp\!\left(2\pi iq\int_L A\right)\to \exp\!\left(-iq\int_{\partial L}\dd\theta\right)\cW_q(L)=\ee^{-iq\theta(x)}\cW_q(L)\,. \tag{3.14}
\end{equation}
于是把两者合在一起，我们就得到一个规范不变的构型
\begin{equation}
  \phi(x)\cW_q(L)\,. \tag{3.15}
\end{equation}
换句话说，我们构造了一个非真局域算符 $\phi$，它生活在电荷为 $q$ 的 Wilson 线 $\cW_q$ 的端点处。

% [Tutor Note: p.30 式 (3.11)--(3.14) 的 $2\pi$ 归一化可自洽核对：由 (3.11) 有 $\int_L A\to\int_L A-\frac{1}{2\pi}\int_L \dd\theta=\int_L A-\frac{\theta(x)}{2\pi}$（用 (3.13)），代入 (3.14) 的 $\exp(2\pi iq\int_L A)$ 即得因子 $\ee^{-iq\theta(x)}$。正因 (3.12) 带 $\ee^{+iq\theta}$，两者相乘才规范不变——这也是 (3.15) 存在的根据。若把 (3.11) 的 $2\pi$ 丢掉，$\phi\cW_q$ 就不变性会被破坏，可作检查点。]

\para{屏蔽。} 类似地，通过插入场 $\phi$ 的幂 $\phi^n$，我们可以构造出把 Wilson 线联系起来的非真局域算符
\begin{equation}
  \cW_p\sim \cW_{p+nq}\,,\qquad \forall\, p\in\Z\,. \tag{3.16}
\end{equation}
这就强加了 Wilson 线算符之间的屏蔽；而 (3.4) 附近定义的未屏蔽线算符等价类之群可以写成
\begin{equation}
  \mathsf{D}_1=\frac{\Z}{q\Z}\equiv\Z_q\,, \tag{3.17}
\end{equation}
其中分子中的 $\Z$ 因子是屏蔽之前的 Wilson 线群，分母中的 $q\Z$ 因子是完全屏蔽的 Wilson 线群。

% [原文 p.31]
\para{电 1-形式对称性。} 按照 3.1 节的一般讨论，群 (3.17) 的 Pontryagin 对偶给出该理论的一个 1-形式对称性群
\begin{equation}
  G^{(1)}=\widehat{\mathsf{D}}_1\cong\Z_q\,. \tag{3.18}
\end{equation}
构成这一 1-形式对称性的拓扑余维-2 算符，是 (2.55) 中所描述算符的一个子集，由下式给出
\begin{equation}
  U_n^{(e)}(\Sigma_{d-2})=\exp\!\left(\frac{2\pi i n}{q}\int_{\Sigma_{d-2}}\star F\right),\qquad n\in\{0,1,\cdots q-1\}\,. \tag{3.19}
\end{equation}
(2.55) 中其他的余维-2 算符不再拓扑，因为它们与电荷为 $q$ 的 Wilson 线 $\cW_q$ 具有非平凡的环绕。

% [Tutor Note: p.31 式 (3.19) 的 $n$ 取值范围 $0,\dots,q-1$ 与 (3.18) 的 $\Z_q$ 严格对应：$\exp(2\pi i n/q)$ 是 $q$ 次单位根，只有 $q$ 个不同取值。核对点：取 $q=2$ 时 $U_n^{(e)}$ 只给出 $\pm 1$ 两个算符，正是 $\Z_2$ 1-形式对称性的生成元与其平方。还要注意 (3.19) 与 (2.55) 的关系：(2.55) 是 $U^{(e)}_g(\Sigma_{d-2})=\exp(i\alpha\oint\star F)$、$g=\ee^{i\alpha}\in \U(1)$；令 $g=\ee^{2\pi i n/q}$ 即得 (3.19)，也就是说 (3.19) 是原 $\U(1)$ 1-形式对称性中"$g$ 为 $q$ 次单位根"的那个 $\Z_q$ 子集——其余 $g$ 所对应的算符因与 $\cW_q$ 环绕非平凡而失去拓扑性。]

% [Tutor Note: p.31 第 2 段原文有笔误 “as they they have a non-trivial linking”，重复了一个 “they”；译文照实保留并在此标明。正确英文应为 “as they have a non-trivial linking with the charge $q$ Wilson line $\cW_q$”。判断依据：同句语法与上下文；该重复不改变语义。]

\begin{examplebox}{例 3.2：把 $\SU(2)$ Higgs 化到 $\U(1)$}
让我们从一个处于某个足够高能标的 $\SU(2)$ 规范理论出发。即便它可能不是 UV 完备的（例如在 $d>4$ 时），我们仍把这一理论称为 UV 理论。所考虑的 $\SU(2)$ 规范理论有一个在 $\SU(2)$ 伴随表示中变换的标量场 $\Phi$。给物质场一个真空期望值（vev）$\langle\Phi\rangle$ 会把 $\SU(2)$ 规范群破缺到它的极大环面
\begin{equation}
  U(1)\subset SU(2)\,, \tag{3.20}
\end{equation}
而在由 $\langle\Phi\rangle$ 设定的标度之下，我们得到一个规范群为 $\U(1)$ 的规范理论。同样，即便它可能不是深 IR，我们仍把这一理论称为 IR 理论。

正如我们后面将要讨论的，UV 的 $\SU(2)$ 理论具有电 1-形式对称性
\begin{equation}
  G^{(1)}=\Z_2\,. \tag{3.21}
\end{equation}
由于 Higgs 机制 (3.20) 所产生的 W 玻色子的存在，它可以与 IR 的 $\U(1)$ 理论的电 1-形式对称性相匹配：该 W 玻色子在 IR 的 $\U(1)$ 规范群下带有电荷
\begin{equation}
  q_\Phi=2 \tag{3.22}
\end{equation}
\end{examplebox}

\begin{examplebox}{例 3.3：把 $\SO(3)$ Higgs 化到 $\U(1)$}
我们可以把 $\SU(2)$ 规范群换成 $\SO(3)$ 规范群，其中
\begin{equation}
  SO(3)=SU(2)/\Z_2\,. \tag{3.23}
\end{equation}
在由 $\langle\Phi\rangle$ 设定的标度之下，$\SO(3)$ 规范群被破缺到它的极大环面
\begin{equation}
  U(1)'\subset SO(3)\,, \tag{3.24}
\end{equation}
% [原文 p.32]
其中 $\U(1)'$ 与 (3.20) 中出现的 $\U(1)$ 群的关系为
\begin{equation}
  U(1)'=\frac{U(1)}{\Z_2}\,. \tag{3.25}
\end{equation}
因此该 W 玻色子的电荷为
\begin{equation}
  q_\Phi=1 \tag{3.26}
\end{equation}
这里的电荷是在 IR 的 $\U(1)'$ 规范群下测量的，于是 IR 的 1-形式对称性群为
\begin{equation}
  G^{(1)}=1\,, \tag{3.27}
\end{equation}
这与 UV 的 $\SO(3)$ 规范理论具有平凡 1-形式对称性这一事实一致，后者将在后面讨论。
\end{examplebox}

% [Tutor Note: p.31--32 例 3.2--3.3 要对照看：(3.22) 的 $q_\Phi=2$ 是 $\SU(2)$ 下伴随表示（自旋 1）在极大环面 $\U(1)$ 上的电荷，故 (3.18)/(3.10) 给出 $\Z_2$ 1-形式对称性；而 (3.26) 的 $q_\Phi=1$ 是同一物理态对 $\U(1)'=\U(1)/\Z_2$ 的电荷。核对方法：$\Z_q$ 对称性要求未屏蔽 Wilson 线步长为 $q$；把 $q$ 从 2 换成 1 后 (3.10) 给出 $\Z_1=1$，即 (3.27)。这也解释了为何 $\SO(3)$ 理论没有电 1-形式对称性。]

\begin{examplebox}{例 3.4：把 $\U(1)$ Higgs 化到 $\Z_q$}
现在我们把上面两个例子的情形反过来：我们不在 IR 中实现 $\U(1)$ 规范理论，而是从一个处于 UV 的 $\U(1)$ 规范理论出发。同时，我们还引入一个电荷为 $q$ 的标量场 $\phi$，并考虑一个具有非零 vev 的真空
\begin{equation}
  \langle\phi\rangle\neq 0\,. \tag{3.28}
\end{equation}
这会使 $\U(1)$ 规范群 Higgs 化到
\begin{equation}
  \Z_q\subset U(1)\,, \tag{3.29}
\end{equation}
因为 $\Z_q$ 保持 $\phi$ 不变。于是我们预期 IR 理论是一个形如
\begin{equation}
  \frac{2\pi}{q}\int a_1\cup\delta b_{d-2} \tag{3.30}
\end{equation}
的 $\Z_q$ 规范理论。

这一点可以相当具体地看出，如下所示。把 $\phi$ 分解为
\begin{equation}
  \phi=\rho \ee^{i\alpha}\,, \tag{3.31}
\end{equation}
并在远 IR 中忽略径向模式、用其 vev
\begin{equation}
  \langle\rho\rangle=v \tag{3.32}
\end{equation}
替代它。于是 Lagrange 量成为
\begin{equation}
  \mathcal{L}=\frac{v^2}{2}(\dd\alpha-qA)\wedge\star(\dd\alpha-qA)\,. \tag{3.33}
\end{equation}
% [原文 p.33]
对圆值标量场 $\alpha$ 取对偶，会引入一个 $(d-2)$-形式场 $B$，其规范对称性为 $B\to B+\dd\Lambda$。用这一对偶场写出的 Lagrange 量就是众所周知的 $BF$ 理论：
\begin{equation}
  \mathcal{L}=\frac{q}{2\pi}B\wedge \dd A\,. \tag{3.34}
\end{equation}
事实上，这就是 (3.30) 中 $\Z_q$ 规范理论的一种连续表述。

正如 2.3 节所讨论的，$\Z_q$ 规范理论具有一个 $\Z_q$ 1-形式对称性，从而与 $\U(1)$ 规范理论的 UV 1-形式对称性相匹配。
\end{examplebox}

% [Tutor Note: p.32 式 (3.30) 原文用的是 $\delta$ 而不是 $\dd$：$a_1\cup\delta b_{d-2}$。核对时请注意次数匹配——$a_1$ 为 1-形式，若 $\delta$ 读作余微分（codifferential）则 $\delta b_{d-2}$ 是 $(d-3)$-形式，$a_1\wedge\delta b_{d-2}$ 只有 $(d-2)$ 次，无法在 $d$ 维上积分；因此这里应把 $\delta$ 读作离散/上链意义下的余边界（coboundary）算子，使 $\delta b_{d-2}$ 为 $(d-1)$-形式，或等价地将其理解为 $\dd b_{d-2}$。此式与 (3.34) 的连续表述 $\frac{q}{2\pi}B\wedge\dd A$ 互为对偶写法，可作交叉验证的锚点。]

% [Tutor Note: p.32--33 例 3.4 给出的是“连续表述”路线：先把 $\phi=\rho \ee^{i\alpha}$ 的径向模式冻到 $v$，得 (3.33) 的 Stückelberg 型 Lagrange 量，再对 $\alpha$ 对偶得到 (3.34) 的 $BF$ 理论。作者说“this is in fact a continuous presentation”一处跳过了论证：$\alpha$ 是周期为 $2\pi$ 的圆值场，其对偶 $B$ 的规范变换 $B\to B+\dd\Lambda$ 保证了 $B$ 的周期性与 $\frac{q}{2\pi}\int B\wedge\dd A$ 的良定义；这一配对只依赖 $q$ 的整性。核对点：把 (3.34) 对 $B$ 变分给出 $\dd A=0$（平坦联络），对 $A$ 变分给出 $\dd B=0$，正是 $\Z_q$ 规范理论（离散 $BF$ 型）的连续极限。]

% [Tutor Note: p.33 式 (3.39) 的商群写法 $\Z/q\Z$ 与 $\Z_q$ 等价；注意 (3.17) 与 (3.39) 是同一个结论在“单物质场”与“多物质场”两种情形下的重述，区别只在 (3.38) 用 $\gcd$ 取代 $q$。这也保证了两处 Pontryagin 对偶 (3.18)、(3.40) 形式一致。]

\para{磁 $(d-3)$-形式对称性。} 即使在加入带电物质场 $\phi$ 之后，磁 $(d-3)$-形式对称性仍保持为
\begin{equation}
  G^{(d-3)}=U(1) \tag{3.35}
\end{equation}
不变。

\subsubsection{含多个物质场的 Maxwell 理论}

现在考虑把多个物质场
\begin{equation}
  \phi_i\,,\qquad 1\le i\le m\,, \tag{3.36}
\end{equation}
加入到纯 Maxwell 理论中，它们的规范电荷为 $q_i$。利用 $\phi_i$ 的乘积，我们可以构造出实现屏蔽的非真局域算符
\begin{equation}
  \cW_p\sim \cW_{p+nq} \tag{3.37}
\end{equation}
其中完全屏蔽掉的最小电荷 $q$，由所有物质场电荷的\textbf{最大公约数（gcd）}给出
\begin{equation}
  q=\gcd(q_1,q_2,\cdots,q_m)\,, \tag{3.38}
\end{equation}
于是未屏蔽 Wilson 线的群为
\begin{equation}
  \mathsf{D}_1=\Z/q\Z\equiv\Z_q\,, \tag{3.39}
\end{equation}
其 Pontryagin 对偶就是电 1-形式对称性群
\begin{equation}
  G^{(1)}=\widehat{\mathsf{D}}_1\cong\Z_q\,. \tag{3.40}
\end{equation}
磁 $(d-3)$-形式对称性则仍然是
\begin{equation}
  G^{(d-3)}=U(1)\,. \tag{3.41}
\end{equation}
